\section*{Chapter \ref{ch:m2:short-term-interest-rate-futures} --- Short-Term Interest-Rate Futures}

\begin{solution}{exo:m2:short-term-interest-rate-futures:1}
$100 - 95.965 = 4.035\%$. A rise of 3.5 basis points is 7 half-basis-point ticks;
on 200 contracts the gain is $3.5 \times 25 \times 200 = \text{USD}~17\,500$.
\end{solution}

\begin{solution}{exo:m2:short-term-interest-rate-futures:2}
17 March, 16 June, 15 September and 15 December 2027. The June 2027 contract's
\omterm{def:m2:short-term-interest-rate-futures:imm}{reference quarter} runs from 17 March (included) to 16 June 2027 (excluded).
\end{solution}

\begin{solution}{exo:m2:short-term-interest-rate-futures:3}
$100 - 96.1025 = 3.8975\%$. A long position loses when rates rise: $2 \times 41.67
\times 100 = \text{USD}~8\,334$.
\end{solution}

\begin{solution}{exo:m2:short-term-interest-rate-futures:4}
Only three of October's 31 days carry the new rate, so the multiplier is $31/3$:
a 1-basis-point price error moves the implied rate by 10.3 basis points, 41\% of a
hike. Use November's contract, in which the post-meeting rate applies every day
(there is no November meeting): its price reads that rate directly.
\end{solution}

\begin{solution}{exo:m2:short-term-interest-rate-futures:5}
$[(1 + 0.0395 \times 91/360)(1 + 0.0405 \times 91/360)(1 + 0.0410 \times 91/360)(1 +
0.0412 \times 91/360) - 1]/(4 \times 91/360) = 4.1178\%$.
\end{solution}

\begin{solution}{exo:m2:short-term-interest-rate-futures:6}
The two exchanges name their contracts differently. The three-month SOFR
``December'' contract compounds from mid-September to mid-December: the
period \emph{ending} in December. The three-month SONIA ``December'' contract
accrues from the third Wednesday of December to mid-March: the period
\emph{starting} in December. The hedge is one quarter out: read the reference
period in each contract's rules, never the month in its name.
\end{solution}

\begin{solution}{exo:m2:short-term-interest-rate-futures:7}
At 1\%: $\frac12 \times 0.0001 \times 5 \times 5.25 = 13.1$ basis points at five years;
$\frac12 \times 0.0001 \times 10 \times 10.25 = 51.25$ at ten. At 1.2\%, both scale by
$1.44$: 18.9 and 73.8 basis points, an increase of 5.8 and 22.55.
\end{solution}

\begin{solution}{exo:m2:short-term-interest-rate-futures:8}
The October contract has only three days at the new rate, so its implied
probability carries ten times the price noise of November's: a difference of
one percentage point of probability is a fraction of a tick in October's
price. And the two numbers are not claims on the same thing: October's
average also contains three days of whatever the month-end turn is. Neither
difference is an arbitrage; the November reading is the reliable one.
\end{solution}

\begin{solution}{pb:m2:short-term-interest-rate-futures:1}
\textbf{1.} $100 - 96.1025 = 3.8975\%$.
\textbf{2.} Nine (1 to 9 December) at the old rate, twenty-two at the new.
\textbf{3.} $(31 \times 3.96 - 9 \times 3.8975)/22 = 3.9856\%$.
\textbf{4.} $(3.9856 - 3.8975)/0.25 = 35.2\%$.
\textbf{5.} One point of probability is 0.25 basis points of the new rate on 22 of
31 days: $0.25 \times 22/31 = 0.18$ basis points of December price.
\textbf{6.} $(31 \times 3.96 - 0.10 - 9 \times 3.8975)/22 = 3.9810\%$: 33.4\%.
\textbf{7.} 3.9765\%: 31.6\%.
\textbf{8.} $0.10/22 \times 100/0.25 = 0.18$ percentage points per basis point of
turn.
\textbf{9.} A one-month contract averages the calendar days of its month; the 31
December fixing covers 31 December, which is in December. January's first days
use the fixing of the first business day of January.
\textbf{10.} From contracts or deposits that straddle the year-end and others that
do not, as in \cref{prop:m2:money-markets:turn}; or from the history of past
year-end fixings as a prior.
\textbf{11.} On a hold the December average is 3.8975\%, the final price 96.1025:
$+6.25$ basis points, $6.25 \times 41.67 \times 1\,000 = +\text{USD}~260\,438$.
\textbf{12.} On a hike the average is $(9 \times 3.8975 + 22 \times 4.1475)/31 =
4.0749\%$, price 95.9251: $-11.49$ basis points, $-\text{USD}~478\,869$.
\textbf{13.} $0.648 \times 260\,438 - 0.352 \times 478\,869 = 0$, to the dollar: the
market's probability is the one at which the trade is fair.
\textbf{14.} At 35.2\%, because the price is exactly the probability-weighted
average of the two settlements; the desk wins on average only if its
probability is lower than the market's.
\textbf{15.} The spread of SOFR to the policy rate (repo conditions, reserves),
month-end and year-end turns, and the realised fixings of the days before the
meeting.
\textbf{16.} The outcomes are not only 0 and 25: a 50-basis-point move or an
off-meeting move are possible; the price mixes all outcomes, so the
two-outcome split is a convention.
\textbf{17.} Futures prices embed risk premia: investors pay to hedge a hike, so
the price-implied probability can exceed the real-world one; and the
Committee's projections are forecasts, not bets.
\textbf{18.} The rate going into December would be about $3.87 + 0.25 = 4.12\%$; the December price
would have to be re-read against it, and markets usually reprice the whole
path after a surprise.
\textbf{19.} Named result: \emph{the implied December hike probability} is 35.2\% with
no turn and 33.4\% with a 10-basis-point turn on 31 December.
\textbf{20.} The average of the overnight rate over a month, and so every
decision and every turn that falls in it.
\end{solution}

\begin{solution}{iq:m2:short-term-interest-rate-futures:1}
Take the contract covering the meeting's month (or better, the month after if
the meeting is late); infer the rate before the meeting from the previous
month's contract; solve the month's average for the post-meeting rate; divide
the change by the size of the move. Remove turns and the SOFR--policy spread
first, and remember it is a risk-neutral, two-outcome approximation.
\iqlookfor{the day-weighting, the late-meeting problem, and the caveats.}
\end{solution}

\begin{solution}{iq:m2:short-term-interest-rate-futures:2}
So that the price moves like a bond price: up when rates fall. Buying the
future is then a long position in rate-sensitive assets, like buying a bill,
and exchanges can list prices, bids and offers in the usual direction; a fixed
value per basis point makes risk the number of contracts.
\iqlookfor{the bond-like convention and the fixed basis-point value.}
\end{solution}

\begin{solution}{iq:m2:short-term-interest-rate-futures:3}
Futures are margined daily, forwards settle once. A short future gains when
rates rise, and reinvests the gains at high rates; it loses when rates fall and
finances the losses at low rates. That correlation benefits the short, so the
futures rate exceeds the forward rate; in a Gaussian model by about
$\frac12\sigma^2 T_1 T_2$, negligible at the front, tens of basis points at ten years.
\iqlookfor{the direction and its reason, and the scaling with horizon.}
\end{solution}

\begin{solution}{iq:m2:short-term-interest-rate-futures:4}
A pack is four consecutive quarterly contracts traded together: one year of
the strip. The second year's pack isolates expectations for rates a year ahead,
the policy path after the next few meetings, without the front contracts'
exposure to the next meeting and to money-market noise; it is the cleanest
futures position on the medium-term path.
\iqlookfor{packs as a way to trade a segment of the curve.}
\end{solution}

\begin{solution}{iq:m2:short-term-interest-rate-futures:5}
Month-ends: turn premia in the fixings, weekends that make the last fixing
cover several days. Late meetings: the post-meeting rate is estimated from a
few days, so noise explodes; switch to the next month's contract. Year-end: a
large turn that, if not removed, reads as a probability of a move. Also: stale
prices in illiquid months, holiday calendars, and the day the new rate takes
effect.
\iqlookfor{concrete failure modes with their fixes.}
\end{solution}

\begin{solution}{iq:m2:short-term-interest-rate-futures:6}
No: it is a risk-neutral probability, containing the premium investors pay to
hedge the outcome they fear; it also assumes two outcomes and a known turn and
spread. To convert: estimate the risk premium (from the history of implied
against realised decisions, or from surveys), compare with your own model's
probability, and trade only the difference net of costs.
\iqlookfor{risk-neutral versus real-world, and how to estimate the gap.}
\end{solution}
